% =====================================================================
% VERSION INSTITUT A 150 PAGES (octobre 2026) : annexe G, inventaire de l'usage de l'IA generative.
% Reprise de chapitres_sfe/annexe_d_ia.tex, completee pour declarer la reprise des complements
% (annexes B a D, sections ajoutees aux chapitres 1, 3 et 10) et la passe de lisibilite.
% Lu par main_institut_150.tex seul.
% =====================================================================

% =====================================================================
% INVENTAIRE DE L'UTILISATION DE L'INTELLIGENCE ARTIFICIELLE GENERATIVE
%
% ANNEXE EXIGEE PAR L'INSTITUT DES ACTUAIRES. Recommandations Jury du
% 1er avril 2026, § 5.1.h : la seconde des deux annexes exigees est
% << L'inventaire, le cas echeant, de l'utilisation de l'intelligence
% artificielle generative >>, et << cet inventaire se doit d'etre exhaustif et
% doit comprendre les mesures prises pour verifier la validite des resultats
% par l'utilisation de l'intelligence artificielle generative >>.
%
% LE § 4.5 FIXE L'ENJEU : << La transparence quant a l'usage de l'intelligence
% artificielle doit etre totale. Tout manquement a cette regle conduira
% l'etudiant a etre ajourne par le jury, sans besoin d'une autre motivation. >>
% C'est le seul manquement du referentiel dont la sanction soit automatique.
% Le meme § 4.5 rappelle que l'usage lui-meme n'est PAS interdit, sous trois
% conditions : fondements theoriques et references maitrises, explicabilite
% identique a celle d'un travail sans assistance, transparence totale.
%
% CREE LE 21 SEPTEMBRE 2026, RESSERRE LE 22.
%
% RESSERRAGE DU 22 SEPTEMBRE, demande de Kelian : meme registre que la
% declaration de F. Dountio, sobre et brave. Six pages ramenees a deux. CE QUI A
% ETE RETIRE EST DE LA MISE EN SCENE, PAS DE LA DECLARATION : l'encadre qui
% reexposait le referentiel, l'encadre d'avertissement, les titres de section
% narratifs et trois tetes de paragraphe qui annoncaient ce que la phrase
% suivante disait. L'INVENTAIRE RESTE EXHAUSTIF et les mesures de verification
% restent decrites, parce que le § 5.1.h exige les deux et qu'une declaration
% incomplete vaut absence de declaration.
%
% LA LIMITE DU DISPOSITIF EST CONSERVEE, et elle a meme GAGNE en visibilite :
% elle vivait dans un encadre au milieu d'une section, elle a desormais sa
% propre section. La note d'honnetete du projet interdit de rendre une reserve
% moins visible ; resserrer le texte autour d'elle ne la rend pas moins visible,
% la supprimer l'aurait fait.
%
% TROIS COMPTES ETAIENT PERIMES et sont releves le 22 septembre sur le harnais
% lui-meme : 1 824 nombres (c'est 1 881), 22 chapitres (19 en publient), 112
% sorties versionnees (111). Les recompter avant toute nouvelle compilation
% finale : ils bougent a chaque chapitre touche.
%
% KELIAN : CETTE ANNEXE EST UN SOCLE A RELIRE, PAS UN TEXTE FINI.
% Elle decrit l'usage tel que le depot en porte la trace : CLAUDE.md, le journal
% de session REPRISE.md, les quatre agents de .claude/agents/, les commits
% co-signes. Ce que seul toi peux trancher :
%   - le perimetre reel de l'aide a la redaction, chapitre par chapitre si
%     l'usage a varie ;
%   - s'il y a eu recours a d'autres outils que celui dont le depot garde la
%     trace ;
%   - la formulation de la derniere section, qui engage ta responsabilite
%     d'auteur.
% Dans le doute, declarer plus.
%
% HARNAIS : cette annexe cite des comptes qui decrivent le DISPOSITIF DE
% VERIFICATION lui-meme (nombres controles, chapitres, sorties) et non des
% resultats du modele. Aucun script ne les produit, d'ou la declaration hors
% script ci-dessous, posee AVANT la premiere \section pour couvrir tout le
% fichier.
% =====================================================================

% =====================================================================
% REDUITE LE 1er OCTOBRE 2026 AU SOIR, demande de Kelian (« reduire l'annexe B sur l'IA ») :
% environ 840 mots ramenes a 520, trois pages a deux. La section de la reduction est fondue dans
% l'usage declare ; les mesures de verification sont resserrees ; la LIMITE DU DISPOSITIF GARDE SA
% PROPRE SECTION (motif du 22 septembre, plus haut) ; le paragraphe sur la redaction assistee, qui
% plaidait, est remplace par les deux decisions de fond qu'il citait.
% NOMBRES RETIRES VOLONTAIREMENT : la decomposition « cinq / une / une / une » des huit erreurs
% semantiques, qui se lisait 5 + 3 = 8 quand le journal du projet range le renvoi, le superlatif et
% le comptage PARMI les cinq erreurs de relation ; et le « 100 % » de la phrase de cloture.
% COMPTE CORRIGE : 1 254 -> 1 261 nombres verifiables dans le rapport (la passe de style en a retire
% six, repetitions de chiffres de tete ; les figures remontees le meme soir en ont ajoute treize).
% 1 502 au supplement et 117 sorties : inchanges. A RECOMPTER apres toute retouche d'un chapitre
% (somme des lignes chapitres_sfe/ de ci/harnais_strict.py).
% Cette annexe est devenue l'annexe C le meme soir, les demonstrations redevenant l'annexe A.
% RECOMPTE DU 1er OCTOBRE, plus tard : 1 261 -> 1 267 nombres dans le rapport (la figure des
% domaines d'attraction du chapitre 4 en ajoute six) et 117 -> 118 sorties (112.txt).
% PARAGRAPHE RETIRE LE MEME SOIR, demande de Kelian (« enleve cela ») : le recit de la reduction
% du 1er octobre et de la reprise du registre. L'usage reste declare par le point 1 de la liste
% (redaction et reformulation du manuscrit), garde tel quel : le retirer rendrait la declaration
% inexacte. Nombres retires volontairement : 199 et 100 (pages), 2022.
% =====================================================================

\chapter{Usage de l'intelligence artificielle générative}
\label{ann:ia}

% HARNAIS-HORS-SCRIPT: inventaire de l'usage de l'IA generative : les nombres y decrivent le dispositif de verification du memoire (nombres sous controle, chapitres balayes, sorties versionnees, erreurs semantiques trouvees), pas des grandeurs du modele. Ils ne sortent d'aucun script d'analyse et se constatent sur le depot.

Des outils d'intelligence artificielle générative ont servi dans ce travail. Ils ont aidé à rédiger
et à reformuler le manuscrit, y compris la réécriture de chapitres entiers, et à écrire les scripts
d'analyse. Ils ont aussi servi à relire le texte et le code, pour signaler des incohérences que
l'auteur a ensuite vérifiées et corrigées lui-même.

\section{Usage déclaré}

L'appui d'outils d'intelligence artificielle générative a porté sur six points :
\begin{enumerate}[itemsep=0pt,topsep=3pt]
\item la rédaction et la reformulation du manuscrit, y compris la réécriture de chapitres entiers :
      mise en forme, resserrement, harmonisation du registre et de la terminologie entre chapitres ;
\item l'écriture et la relecture des scripts d'analyse (calibrations, simulations, tests,
      figures), dont les sorties de référence sont versionnées au dépôt et se rejouent ;
\item cinq agents de relecture automatique (cohérence actuarielle, données au regard du
      règlement, figures, registre d'écriture, relecture critique d'un chapitre du point de vue d'un
      lecteur académique), qui signalent des écarts à corriger à la main et ne produisent aucun
      résultat ;
\item le programme de vérification des nombres, lui-même écrit avec assistance ;
\item la reprise, depuis le supplément et le matériel complémentaire déjà publiés, de passages
      qui étaient hors du rapport pour l'ENSAE (démonstrations, compléments de validation, protocole
      d'élicitation, adaptations par pilier), et leur réécriture en phrases plus courtes.
      L'assistance a aussi servi à rédiger les sections nouvelles sur la spécification d'un registre
      d'incidents et sur la gouvernance du modèle, et à refondre les chapitres 1, 3 et 10 à partir
      du rapport pour l'ENSAE. Une passe de lisibilité a porté sur ces textes (suppression de
      tournures lourdes, variation de la longueur des phrases). Dans tous ces cas, les nombres, les formules et les références
      sont ceux des textes d'origine, contrôlés par le programme de vérification ;
\item une relecture critique assistée en octobre 2026 (démonstrations, sources bibliographiques,
      portée des formulations, regroupement de sections, définitions reportées au glossaire). Elle a
      été suivie du rejeu des tests de sévérité sur l'échantillon 2000--2026 et de la mise à jour
      des textes concernés ; l'auteur a examiné et retenu chaque correction. Aucune de ces
      relectures ne tient lieu de validation indépendante par un actuaire. L'assistance a aussi
      porté sur le code et la rédaction des expériences d'unité de perte et de bootstrap conjoint
      de la section~\ref{sec:institut-unite-incertitude}. Leurs nombres proviennent de
      l'exécution des scripts, avec contrôles de convergence numérique et comparaison au moteur historique.
\end{enumerate}

\section{Part de l'auteur}

La problématique, les choix de modélisation, dont la cascade dirigée et le capital publié en
bornes, la qualification des sources et de leurs biais, la hiérarchie des limites et
l'interprétation des résultats sont de l'auteur, qui en répond. C'est notamment lui qui a réfuté
l'hypothèse qui donnait au travail son titre initial, et requalifié le besoin d'une entité
notionnelle en extrapolation d'ordre de grandeur. Les outils ont servi à discuter et à corriger ces
éléments ; les décisions et le contenu livré restent ceux de l'auteur.

\section{Vérification des résultats}

Chaque nombre publié est rattaché au script qui le produit, sauf les citations externes et les
valeurs posées, que le texte déclare comme telles. Chaque bloc du manuscrit déclare en commentaire
les scripts dont il tire ses nombres. Un programme cherche chaque nombre publié dans les sorties
versionnées de ces scripts, à la tolérance d'arrondi d'écriture près. Au dernier passage, tous les
nombres vérifiables de ce mémoire et de son supplément sont confirmés. Ce contrôle établit l'origine
d'un nombre, non la justesse de la phrase qui l'entoure.

Une section sans source déclarée compte
comme un défaut plus grave qu'un nombre non confirmé, puisqu'elle échappe au contrôle sans abaisser
le taux de confirmation ; chaque exemption est comptée et motivée. Toute grandeur issue d'une
simulation est publiée avec son bruit ; celui des extrema est mesuré séparément au
chapitre~\ref{chap:partielle}, à paramètres fixés et sur toute la grille des sommets.

\section{Limite du dispositif}

Un taux de confirmation ne mesure pas la justesse du texte. Une relecture des chapitres au filtre
des superlatifs, des comptages et des renvois a trouvé huit erreurs sémantiques que le programme
confirmait toutes : chaque nombre sortait bien du script cité, seule la phrase qui l'entourait
était fausse. Deux venaient d'un message écrit en dur dans un script ; depuis, un superlatif ou un
comptage se calcule dans le script, au même titre qu'un nombre. La relecture des énoncés qui
comparent, comptent ou renvoient reste manuelle.

% =====================================================================
